\documentclass{article}

% Language setting
% Replace `english' with e.g. `spanish' to change the document language
\usepackage[english]{babel}

% Set page size and margins
% Replace `letterpaper' with `a4paper' for UK/EU standard size
\usepackage[letterpaper,top=2cm,bottom=2cm,left=3cm,right=3cm,marginparwidth=1.75cm]{geometry}

% Useful packages
\usepackage{amsmath}
\usepackage{graphicx}
\usepackage[colorlinks=true, allcolors=blue]{hyperref}


\title{Answering LE 105 Assignment 3}
\author{Chaitanya Deshkar | 23B2712}

\begin{document}

\maketitle

\section*{Question 1}
\subsection*{Part a}
The reason why a solar eclipse occurs, is because the moon comes in between the sun and the earth, and thus, the light from the sun is blocked by the moon. This can only happen when the sun, moon and earth are collinear. 
\newline
While the reason why lunar eclipse occurs is because the earth comes in between the sun and the moon, and thus, the light from the sun is blocked by the earth. This can only happen when the sun, moon and earth are collinear. This happens because
moon doesn't have the light of its own, and it only reflects the light from the sun. So, when the earth comes in between the sun and the moon, the light from the sun is blocked by the earth, and thus, the moon doesn't get any light from the sun, making the moon disappear.
\subsection*{Part b}
Solar Eclipse occurs on the amawasya tithi, or the new moon day of the month, when the moon comes in between the sun and the earth. Lunar Eclipse occurs on the poornima (full moon) day of the month, when the earth comes in between the sun and the moon.
\subsection*{Part c}
The reason why eclipses don't occur every month is because the moon's orbit is tilted by about 5 degrees to the earth's orbit around the sun. So, most of the time, the moon is either above or below the line between the sun and the earth, and thus, eclipses don't occur every month.
\subsection*{Part d}
This is just subject to the sizes of the sun, moon and earth, and their distances from each other. The sun is about 400 times larger than the moon, but it is also about 400 times farther away from the earth than the moon. So, they both appear to be about the same size in the sky. The earth is about 3.67 times larger than the moon, so during a lunar eclipse, the earth's shadow completely covers the moon. But because the moon is smaller than sun, the duration of solar eclipse is much shorter than that of lunar eclipse.
\newline

\section*{Question 2}
The given equation reduces to:
\[354.3672n + 29.5306m = 365.2594n\]
\[\Rightarrow 10.8922n = 29.5306m\]
\[\Rightarrow \frac{n}{m} = \frac{29.5306}{10.8922}\]
meaning nearly for every 30 lunar months, there are 11 solar months. This is the reason why there are 7 leap months in a 19 year cycle, because in 19 years, there are nearly 235 lunar months and 228 solar months. So, to keep the lunar calendar in sync with the solar calendar, we need to add 7 leap months in a 19 year cycle.

\end{document}